\documentclass[11pt]{article}
\usepackage[T1]{fontenc}
\usepackage[margin=0.68in]{geometry}
\usepackage{amsmath,amssymb,booktabs,array}
\usepackage[hidelinks]{hyperref}
\setlength{\parindent}{0pt}
\setlength{\parskip}{3pt}
\emergencystretch=2em
\renewcommand{\arraystretch}{1.08}
\urlstyle{same}
\title{\vspace{-1.5em}Selection experiments at order thirteen\\
\large Finite packet progress and restricted screens; no completed lift claim}
\author{Theo H.}
\date{22 November 2026}
\begin{document}
\maketitle
\vspace{-1.5em}
This note records selection-stage experiments and restricted RLLB screens.
\textbf{It claims neither a new N=13 record nor a completed recursive lift.}
The supplied archive contains native selections, local motifs, and previous
witnesses with differing replay status. They are not interchangeable forms
of evidence for a completed new word.

\section*{Objects and conventions}
Use Pantone's \emph{transport} only for its defined full/short-row map~\cite{pantone}.
At $n=13$, $K=11$, a complete native selection chooses one row or no row in
$10!=3{,}628{,}800$ 2-loops. Completion concerns
$12!=479{,}001{,}600$ cyclic classes (1-cycles) on thirteen symbols;
this is not the number $13!$ of individual permutations.
A packet is a complete balanced solution on a residual $(7,4)$ block of
5,040 2-loops. The port-monodromy product counts its completed trails.
Terminology follows~\cite{glossary,spra}.

\section*{1. Why a monodromy-correct motif was insufficient}
The original 108-step itinerary had port cycles $(3,7)$ but visited only
99 distinct native 2-loops, repeating nine. Removing two identity-action
circuits preserved the port action and left an 88-step itinerary with one
repeat. A tested zero-repeat 98-row motif had charge 232, but remained a
partial walk until integrated into a full balanced selection.

Conservative integration demoted entire intersected baseline walks to
loops without rows. It guaranteed closure but could produce too many
completed trails. The best recorded compound integration has
$Q=1{,}786{,}938$, $T=81{,}970$. Its intact weight-3 floor is
\[
F_3(13)+Q+T+9=6{,}747{,}808{,}117,
\]
which misses the experiment's $6{,}747{,}800{,}000$ target.
That is a diagnosis of this integration and model, not an impossibility
proof against nonlocal repairs, other openings, or a future lift.

\section*{2. Balanced switches yielded a complete finite packet}
The Step1-to-N packet difference contains exactly 259 changed 2-loops,
partitioned into eight independent balanced switches: one component of
161 loops and seven components of 14. Enumerating all $2^8=256$ switch
combinations gave 56 with minimum $Q+T=2,869$ in this finite family.

A 42-row patch produces a complete packet with $Q=2,443$, 406 loops without
a row, five active row walks, 20 active completion trails, and $T=426$.
One 288-row walk has monodromy cycles $(7,3)$. This realizes an A-like
finite port pattern at this degree; it does not prove that the original
N=11 packet interiors possess an ordinary recursive transport.

Applying the packet in 48 residual fibers gives the archived complete
native totals
\[
Q=1{,}812{,}288,\qquad D=19{,}488,\qquad T=28{,}164.
\]
Original native and completion-class receipts are retained as archived
evidence. No new 6.7-billion-symbol literal word was assembled from this
selection in this repository session.

\section*{3. What the screens and mask comparisons establish}
For one intact weight-3 opening per completed trail, $h=10$, so
\[
L=F_3+Q+10T-O\geq F_3+Q+T+9=6{,}747{,}779{,}661.
\]
The archive's grouped global trie gave $6{,}747{,}779{,}651$, ten weaker.
A lower bound below the comparison word is only a necessary screen;
it is not an achievable arrangement cost.

Mask 186 changes 2,688 native row assignments relative to mask 38 across
48 fibers and leaves the displayed $Q,D,T$ totals unchanged. All 56
minimum-floor variants have the same simple relaxed matching counts.
Seed-dependent local phase-compatible path covers motivated mask 186 as
a follow-up; they do not establish a globally better word or an optimal mask.

A separate selection-first compound-core checkpoint lowers $Q+T$ by 780
against its retained N-block control. Replacing all residual blocks with
Step1 gives a structural gain of 1,836 against that experiment's Halfar-N
baseline. These counts must not be reported as literal-length savings.

\section*{4. Current computational chain and next obligation}
The first chain in the companion guide builds and runs
\path{search/rllb.cpp} as \path{build/rllb} on a complete candidate's
port graph. It offers the archived grouped endpoint-symbol price family
or full connector keys; they are different relaxations. Both preserve one
shared opening phase per trail, and neither certifies unrestricted splits
or weight-2 openings.

The exact progress is a finite balanced packet, complete experimental
native selections, and a clearer accounting of why some integrations fail.
The remaining obligation is a constructive, connection-aware global
arrangement with a separately identified literal witness. Until that
obligation is met, the correct status is selection research and attempted
extension, not a completed lift or record.

\section*{Sources and reproducibility}
The companion \href{REPRODUCE.md}{\texttt{REPRODUCE.md}} gives the current
RLLB-first chain, pinned inputs, and exact replay scope;
\href{NOTES.md}{\texttt{NOTES.md}} separates personal research notes.
The source package and retained receipts are identified by
\href{source-import.json}{\texttt{source-import.json}}.

\begin{thebibliography}{9}\footnotesize
\setlength{\itemsep}{2pt}
\bibitem{pantone} J. Pantone,
\emph{Superpermutations of length
$n!+(n-1)!+(n-2)!+(43/80+o(1))(n-3)!$}, brief summary,
September 2026, updated October 2026, Sections 2--7.
\url{https://github.com/jaypantone/superperm-upper-43-80/blob/main/docs/upper-bound-explained.pdf}.
\bibitem{glossary} J. Halfar, \emph{Superpermutation Glossary},
21 November 2026 evening version; supplied 35-page printout,
\href{../../rllb/sources/Glossary.pdf}{archived source}, especially pp.~14--20 and 32--35.
\bibitem{spra} T. H., \emph{The Reduced Alphabet: Encoding Superpermutations
by Transport Events}, SPRA 0.2 reconciliation, 21 November 2026.
\href{../../rllb/sources/the-reduced-alphabet-spra-0.2-reconciled.md}{Supplied note}.
\bibitem{almanac} \emph{Superpermutation Almanac}, record notes and word identities,
\href{../../rllb/almanac-snapshot.json}{snapshot of 22 November 2026}.
\url{https://superpermutations.com/api/words}.
\end{thebibliography}
\end{document}
